
\subsection{Notation}

\paragraph{Time and baseline variables.}
Let $t = 0, 1, 2, \ldots$ index days since Call Center contact. Baseline covariates measured at $t = 0$ are denoted by $\boldsymbol{L}$.

\paragraph{Treatment.}
Let $T_A = \min_s \{s : A(s) = 1\}$ denote the time of SSVF initiation. Define $A(t) = \mathbf{1}[T_A \leq t]$ as the indicator of having initiated SSVF at or by day $t$; once $A(t) = 1$, it remains 1 (absorbing). The grace period treatment indicator is $\Delta_A = \mathbf{1}[T_A \leq 30]$.

\paragraph{Housing status.}
A patient's housing status $H(t)$ is binary: $H(t) = 1$ if housed and $H(t) = 0$ if homeless. $H(t)$ is only defined for living individuals, i.e., when $D(t) = 0$.

\paragraph{Mortality.}
Let $D(t) = 1$ if the patient has died by day $t$ and $D(t) = 0$ otherwise. Mortality is always observed. When $D(t) = 1$, $H(t)$ is undefined. Let $T^D$ denote the time of death.

\paragraph{SSVF exit.}
Let $Z(t) = 1$ if the patient has initiated and exited SSVF by day $t$, and $Z(t) = 0$ otherwise. $Z(t)$ is binary and absorbing; $Z(t) = 1$ implies $A(t) = 1$. Let $T^Z$ denote the day of SSVF exit. The exit destination is captured by $H(T^Z)$, which is assumed observed.

\paragraph{Time-varying covariates.}
Let $\boldsymbol{X}(t)$ represent time-varying covariates:
\begin{itemize}
    \item $X_1(t)$: Indicator that the Veteran is enrolled in HUD-VASH
    \item $X_2(t)$: Indicator that the Veteran is enrolled in HUD-VASH and living in HUD-VASH housing
    \item $X_3(t)$: Indicator that the Veteran is enrolled in GPD or other VA homeless programs
    \item $X_4(t)$: The individual's VA enrollment status
\end{itemize}

\paragraph{Healthcare visits and data recorded on day $t$.}
Let $V(t)$ represent the visit type at time $t$:
\begin{itemize}
    \item $V(t) = 0$: No visit
    \item $V(t) = 1$: Homeless clinic visit
    \item $V(t) = 2$: HUD-VASH visit
    \item $V(t) = 3$: Risk visit
    \item $V(t) = 4$: Mental health or substance use visit
    \item $V(t) = 5$: Social work visit
    \item $V(t) \geq 6$: Other visits (primary care, emergency, etc.)
\end{itemize}

Define $B(t) = \{D(t),\, A(t),\, Z(t),\, \boldsymbol{X}(t),\, V(t),\, \boldsymbol{M}(t)\}$ as the data recorded on day $t$.

\paragraph{Manifest variables.}
When $V(t) \in \{1, \ldots, 5\}$, a vector of manifest variables $\boldsymbol{M}(t)$ is observed. For $V(t) = 0$ or $V(t) \geq 6$, $\boldsymbol{M}(t) = \emptyset$. The manifest variables contain noisy indicators of housing status extracted from ICD-10 codes and clinical notes:
\begin{itemize}
    \item $I(t)$: ICD-10 housing codes. In settings where granular coding is available, $I^G(t) \in \{\text{absent},\, \text{risk},\, \text{homeless},\, \text{sheltered},\, \text{unsheltered}\}$. Otherwise, $I^C(t) \in \{\text{absent},\, \text{present}\}$.
    \item $N(t)$: Count of clinical notes at the visit.
    \item $N_{\text{stable}}(t)$, $N_{\text{unstable}}(t)$, $N_{\text{unknown}}(t)$: Notes classified as stably housed, unstably housed, and unknown, respectively, with $N_{\text{stable}}(t) + N_{\text{unstable}}(t) + N_{\text{unknown}}(t) = N(t)$.
    \item $N_{\text{homeless}}(t)$, $N_{\text{housed}}(t)$, $N_{\text{need}}(t)$: Counts of homeless, housed, and housing need entity spans identified across clinical notes by NLP.
\end{itemize}

Housing status is also directly observed on certain occasions through the structural components of $B(t)$: at baseline ($t = 0$) through the Call Center assessment, at SSVF entry ($t = T_A$), at SSVF exit ($t = T^Z$), when $X_3(t) = 1$ (GPD enrollment implies $H(t) = 0$), and when $X_2(t) = 1$ (HUD-VASH housing implies $H(t) = 1$).

\paragraph{History notation.}
For an arbitrary process $W$:
\begin{itemize}
    \item $\overline{W}(t) = \{W(u) : 0 \leq u \leq t\}$: history through time $t$ (inclusive)
    \item $\overline{W}(t-) = \{W(u) : 0 \leq u < t\}$: history prior to time $t$ (exclusive)
\end{itemize}

\subsection{Outcome, potential outcomes, and estimand}

\paragraph{Composite outcome.}
The primary outcome at time $t$ is $Y(t) = (1 - D(t)) \cdot H(t)$, an indicator of being alive and housed. The follow-up time is $\tau = T_A + 120$.

\paragraph{Potential outcomes.}
For a subject alive at time $s$, define potential outcomes under treatment regime $a \in \{0, 1\}$:
\begin{itemize}
    \item $D^a(t, s)$: Indicator of having died by day $t$
    \item $B^a(t, s)$: Data recorded on day $t$
    \item $H^a(t, s)$: Housing status at day $t$ (defined only if $D^a(t, s) = 0$)
    \item $Y^a(t, s) = (1 - D^a(t, s)) \cdot H^a(t, s)$: Composite outcome at day $t$
\end{itemize}
Here $a = 1$ denotes SSVF initiation at time $s$ and $a = 0$ denotes no SSVF. The index $s$ encodes that the patient is alive and eligible at day $s$ and sets the time origin for follow-up. Under consistency, $Y^1(t, T_A) = Y(t)$ for treated patients.

\paragraph{Estimand.}
The estimand is the average treatment effect among the treated (ATT):
\begin{align*}
    \delta = E\!\left[Y^1(T_A + 120,\, T_A) - Y^0(T_A + 120,\, T_A) \mid \Delta_A = 1\right]
\end{align*}
